\noindent
Static application security testing (SAST) applies static program analysis to the security of software projects. Static analysis is performed on the source code, compiled byte-code, or some other representation of the source code \cite{principles_program_analysis}. Popular SAST tools include: SonarQube, SpotBugs, Snyk Code, and Facebook's Infer. SAST tools have become widely used at all stages of the software development life-cycle. Furthermore, they have been demonstrated to be highly effective in identifying defects in code compared to other methods, such as manual inspection \cite{costs_static_analysis}.\\

\noindent
Unfortunately, SAST tools are not perfect. From Rice's Theorem, it is known that any non-trivial property about programs is reducible to the halting problem, i.e., undecidable \cite{moller2012static}. That is, a perfect analysis is not always possible, and approximations must be made in order for the analysis to terminate, resulting in an imperfect output. False positives (where safe code is incorrectly labelled as a defect) and false negatives (where the analysis misses a genuine fault) \cite{chess2004static} are common in all static analysis tools. False positives are not ideal as manual inspection is required to isolate genuine defects, leading to an increased time and labour cost. And, false negatives of course mean that genuine software faults are left undetected by the analysis. While solving the halting problem may be out of scope for an undergraduate thesis, there are still additional methodologies for improving existing static analyses that do not involve upsetting Mr. Turing.\\

\noindent
To reduce the rate of false negatives, manual configuration of the SAST tools for specific software projects is often required \cite{ZampettiFiorella2017HOSP}. For example, to perform an effective taint analysis, a custom configuration of sources, sinks and sanitisers for a particular project is often required. These taint analysis tools are not smart enough to know if a method is a source, sink, or a sanitiser --- they require sensitive methods to be pre-labelled to conduct the analysis. Generally, these tools come with a default configuration containing a list of sensitive methods in popular APIs. For example, the getParamater() method as a part of the ServletRequest interface would likely be listed by all taint analysis tools as a source. Of course, if a project is using a niche library or code written in-house, the taint analysis will not be able to incorporate those data-points into its analysis, leading to a sub-optimal analysis. To improve the taint analysis, one can always identify the sensitive methods manually and add these to the taint analysises configuration. However, creating this manual configuration is often a time-consuming and resource-intensive process. As such, methodologies to automatically label security sensitive methods may enable existing taint analysis tools to be configured more efficiently and effectively.\\

\noindent
Deep learning, a subtopic within the broader machine learning, has recently increased in popularity. What was once a niche topic reserved for academics with access to computational resources far out of reach by most has become arguably one of the most influential technologies of the modern age. Deep learning has fundamentally changed our everyday life, with applications ranging from targeted advertising to identifying COVID-19 cases from lung CT scans \cite{RAHIMZADEH2021102588}. Natural-language processing (NLP) is one often overlooked application of deep learning, but it has seen tremendous growth in recent years.\\